\documentclass[aspectratio=169]{beamer}

% Tema UFS — Universidade Federal de Sergipe
\usetheme{UFS}

\usepackage[utf8]{inputenc}
\usepackage[portuguese]{babel}
\usepackage{amsmath,amssymb,graphicx,booktabs,xcolor}
\usepackage[nolinks]{qrcode}
\usepackage{tikz}
\usetikzlibrary{shapes.geometric,arrows,arrows.meta,positioning,matrix,fit,calc,
  backgrounds,decorations.pathreplacing,shadows,patterns,
  automata,intersections,shapes}

% Listagens — paleta VS Code Light+
\usepackage{listings}
\definecolor{bgcolor}{HTML}{FFFFFF}
\definecolor{linecolor}{HTML}{DDDDDD}
\definecolor{numcolor}{HTML}{999999}
\definecolor{kwcolor}{HTML}{0000FF}
\definecolor{strcolor}{HTML}{A31515}
\definecolor{cmtcolor}{HTML}{008000}

\lstdefinestyle{ufsStyle}{
  language=Python,
  backgroundcolor=\color{bgcolor},
  basicstyle=\ttfamily\scriptsize\color{black},
  keywordstyle=\color{kwcolor}\bfseries,
  stringstyle=\color{strcolor},
  commentstyle=\color{cmtcolor}\itshape,
  numberstyle=\tiny\color{numcolor},
  numbers=left, numbersep=8pt,
  xleftmargin=16pt, framexleftmargin=16pt,
  breaklines=true, tabsize=4,
  keepspaces=true, showspaces=false,
  showstringspaces=false, showtabs=false,
  frame=single, rulecolor=\color{linecolor},
  framesep=3pt, captionpos=b, upquote=true,
  columns=fullflexible,
}
\lstset{style=ufsStyle}

\title{Refatoração em Python}
\subtitle{Cheiros, Refatorações e a Disciplina do Ciclo Red--Green--Refactor}
\author{Prof.\ Dr.\ André Yoshiaki Kashiwabara}
\institute{DCOMP --- Universidade Federal de Sergipe\\
\texttt{andreyoshiaki@dcomp.ufs.br}}
\date{}

\begin{document}

\begin{frame}[plain]
  \titlepage
\end{frame}

\begin{frame}{A narrativa: Problema $\to$ Solução $\to$ Disciplina}
  \begin{center}
  \begin{tikzpicture}[
    item/.style={rectangle, draw=UFSBlue!60, fill=UFSBlue!10, rounded corners,
      minimum width=9cm, minimum height=0.8cm, font=\small},
    arrow/.style={->, thick, UFSBlue}
  ]
    \node[item] (t1) at (0, 1.5)  {Problema: \emph{enxergar} o que está ruim --- cheiros de código};
    \node[item] (t2) at (0, 0.0)  {Solução: \emph{consertar} com refatorações nomeadas};
    \node[item] (t3) at (0,-1.5)  {Disciplina: \emph{sustentar} com Red--Green--Refactor, testes e PEP 8};
    \draw[arrow] (t1) -- (t2);
    \draw[arrow] (t2) -- (t3);
  \end{tikzpicture}
  \end{center}
\end{frame}

% =====================================================================
\part{Problema: Maus Cheiros de Código}
% =====================================================================
\section{O que são \emph{code smells}}

\begin{frame}{Por que falar de \emph{cheiros}?}
  \begin{columns}[T]
    \begin{column}{0.5\textwidth}
      \textbf{Problema:}\\[0.4em]
      Código \emph{funciona} mas é difícil de manter,
      estender ou entender. Bugs aparecem em cascata,
      mudanças simples viram epopeias.
    \end{column}
    \begin{column}{0.5\textwidth}
      \textbf{Solução:}\\[0.4em]
      Aprender a \textbf{farejar} sinais comuns de
      degradação --- e ter uma refatoração específica
      pronta para cada um.
    \end{column}
  \end{columns}
\end{frame}

\begin{frame}{Code Smell: definição}
  \begin{block}{Definição (Kent Beck \& Martin Fowler)}
    Um \emph{code smell} é uma estrutura no código que \textbf{sugere}
    --- mas não prova --- a necessidade de refatoração.
    Não é bug: o código funciona, mas algo está \emph{cheirando mal}.
  \end{block}
  \begin{exampleblock}{Intuição --- ``If it stinks, change it''}
    O nariz é treinável. Quanto mais código você lê,
    mais rápido percebe que algo está errado --- antes
    mesmo de saber explicar o quê. Smells são \emph{atalhos heurísticos}.
  \end{exampleblock}
\end{frame}

\begin{frame}{Smell vs.\ Bug}
  \begin{center}
  \begin{tikzpicture}[
    box/.style={rectangle, rounded corners, draw, minimum width=5.5cm,
                minimum height=2.4cm, align=center, font=\small}
  ]
    \node[box, fill=UFSRed!20]    (bug)   at (-3.5, 0)
      {\textbf{Bug}\\[0.3em]
       Código se comporta errado.\\
       Testes falham (ou deveriam).\\
       \emph{Urgência:} alta.};
    \node[box, fill=UFSYellow!40] (smell) at ( 3.5, 0)
      {\textbf{Smell}\\[0.3em]
       Código se comporta certo,\\
       mas é frágil para mudar.\\
       \emph{Urgência:} média/baixa.};
  \end{tikzpicture}
  \end{center}
  \vspace{0.4em}
  \footnotesize Smells acumulados \textbf{viram} bugs ao longo do tempo --- é dívida técnica vencendo juros.
\end{frame}

\section{As quatro famílias de Fowler}

\begin{frame}{O mapa dos Code Smells}
  \begin{center}
  \begin{tikzpicture}[
    cat/.style={rectangle, draw, rounded corners, minimum width=3cm,
                minimum height=0.7cm, font=\small\bfseries, align=center},
    item/.style={rectangle, draw, rounded corners, fill=white,
                 minimum width=3cm, minimum height=0.55cm, font=\scriptsize, align=center}
  ]
    \node[cat, fill=UFSRed!30]    (a) at (0,3)   {Bloaters};
    \node[cat, fill=UFSYellow!50] (b) at (3.5,3) {OO Abusers};
    \node[cat, fill=UFSBlue!25]   (c) at (7,3)   {Change Preventers};
    \node[cat, fill=green!30]     (d) at (10.5,3){Dispensables};

    \node[item] at (0,2.2)    {Long Method};
    \node[item] at (0,1.5)    {Large Class};
    \node[item] at (0,0.8)    {Long Parameter List};

    \node[item] at (3.5,2.2)  {Switch Statements};
    \node[item] at (3.5,1.5)  {Temporary Field};
    \node[item] at (3.5,0.8)  {Refused Bequest};

    \node[item] at (7,2.2)    {Divergent Change};
    \node[item] at (7,1.5)    {Shotgun Surgery};
    \node[item] at (7,0.8)    {Feature Envy};

    \node[item] at (10.5,2.2) {Duplicated Code};
    \node[item] at (10.5,1.5) {Dead Code};
    \node[item] at (10.5,0.8) {Speculative Generality};
  \end{tikzpicture}
  \end{center}
  \vspace{0.4em}
  \footnotesize Veremos um representante de cada família + os clássicos.
\end{frame}

\begin{frame}{As quatro famílias em uma frase}
  \begin{itemize}
    \item \textbf{Bloaters} --- coisas que \emph{incham} (métodos, classes, listas).
    \item \textbf{OO Abusers} --- recursos de OO usados de forma errada (ou não usados).
    \item \textbf{Change Preventers} --- estruturas que \emph{dificultam} mudanças futuras.
    \item \textbf{Dispensables} --- coisas que \emph{não precisariam estar lá} (duplicação, código morto, abstração à toa).
  \end{itemize}
\end{frame}

\section{Bloaters}

\begin{frame}[fragile]{Long Method (Método Longo)}
  \begin{alertblock}{Sintomas}
    Mais de $\sim$20 linhas, muitos níveis de indentação,
    comentários explicando \emph{blocos} de código.
  \end{alertblock}
\begin{lstlisting}
def processar_pedido(p):
    if not p.cliente: raise ValueError("sem cliente")
    if not p.itens:   raise ValueError("sem itens")
    total = 0
    for i in p.itens: total += i.preco * i.qtd
    if total > 1000: total *= 0.9
    msg = f"Pedido {p.id}: R$ {total}"
    smtp.send(p.cliente.email, msg)
    return total
\end{lstlisting}
  \footnotesize\textbf{Refatoração de combate:} \emph{Extract Function} --- veremos a versão depois na Parte II.
\end{frame}

\begin{frame}{Large Class (Classe Grande)}
  \begin{center}
  \begin{tikzpicture}[
    big/.style={rectangle, draw, fill=UFSRed!25, minimum width=3.5cm,
                minimum height=3.5cm, font=\small, align=center},
    small/.style={rectangle, draw, fill=green!25, minimum width=2.3cm,
                  minimum height=0.55cm, font=\scriptsize, align=center}
  ]
    \node[big] (P) at (0,0) {\textbf{Pedido}\\[0.3em]
      cliente, itens,\\
      total, desconto,\\
      frete, endereço,\\
      pagamento, status\\[0.4em]
      \emph{15 métodos}};

    \node[small] at (5, 1.5)  {Pedido};
    \node[small] at (5, 0.75) {Calculadora};
    \node[small] at (5, 0.0)  {Frete};
    \node[small] at (5,-0.75) {Pagamento};
    \node[small] at (5,-1.5)  {Notificador};

    \draw[->, thick, UFSBlue] (1.9, 0) -- (3.8, 0);
  \end{tikzpicture}
  \end{center}
  \footnotesize\textbf{Refatoração:} \emph{Extract Class} --- separar por responsabilidade (SRP).
\end{frame}

\begin{frame}[fragile]{Long Parameter List}
\begin{lstlisting}
def criar_pedido(nome, cpf, email, rua, numero,
                 cidade, cep, produto, qtd, preco,
                 cupom, forma_pagto):
    ...
\end{lstlisting}
  \footnotesize
  \textbf{Sintoma:} mais de 3--4 parâmetros, especialmente quando vários
  aparecem juntos (\emph{data clumps}).\\
  \textbf{Refatoração:} \emph{Introduce Parameter Object} ---
  agrupar em \texttt{Cliente}, \texttt{Endereco}, \texttt{ItemPedido}, \texttt{Pagamento}.
\end{frame}

\section{Dispensables, Change Preventers e OO Abusers}

\begin{frame}[fragile]{Duplicated Code}
  \begin{alertblock}{O problema}
    A mesma estrutura aparece em mais de um lugar.
    Mudança em um exige caçar e mudar os outros.
  \end{alertblock}
\begin{lstlisting}
def preco_a(x): return x if x > 0 else -x
def preco_b(y): return y if y > 0 else -y
def preco_c(z): return z if z > 0 else -z
# depois: def preco(v): return abs(v)
\end{lstlisting}
  \footnotesize\textbf{Refatorações de combate:} \emph{Extract Function},
  \emph{Pull Up Method} (para superclasse), \emph{Extract Class} (utilitária).
\end{frame}

\begin{frame}{Divergent Change vs.\ Shotgun Surgery}
  \begin{columns}[T]
    \begin{column}{0.5\textwidth}
      \centering\textbf{Divergent Change}\\
      \begin{tikzpicture}[
        cause/.style={ellipse, draw, fill=UFSYellow!50, font=\scriptsize, minimum width=1.8cm},
        target/.style={rectangle, draw, fill=UFSRed!25, font=\scriptsize, minimum width=1.8cm, minimum height=0.8cm, align=center}
      ]
        \node[cause] (c1) at (-1.8, 1.0)  {layout};
        \node[cause] (c2) at ( 1.8, 1.0)  {cálculo};
        \node[cause] (c3) at ( 0,  -1.0)  {banco};
        \node[target] (T) at (0, 0) {\textbf{Relatório}};
        \draw[->] (c1) -- (T); \draw[->] (c2) -- (T); \draw[->] (c3) -- (T);
      \end{tikzpicture}\\
      \footnotesize Mesma classe muda por razões diferentes (viola SRP).\\
      \textbf{Refatoração:} \emph{Extract Class}.
    \end{column}
    \begin{column}{0.5\textwidth}
      \centering\textbf{Shotgun Surgery}\\
      \begin{tikzpicture}[
        src/.style={circle, draw, fill=UFSRed!30, minimum size=0.9cm, font=\scriptsize, align=center},
        tgt/.style={rectangle, draw, fill=UFSBlue!15, minimum width=1.0cm, minimum height=0.5cm, font=\tiny}
      ]
        \node[src] (M) at (-2.4, 0) {fmt};
        \node[tgt] (a) at (-0.6, 1.0) {A};
        \node[tgt] (b) at ( 0.6, 1.0) {B};
        \node[tgt] (c) at ( 1.8, 1.0) {C};
        \node[tgt] (d) at (-0.6,-1.0) {D};
        \node[tgt] (e) at ( 0.6,-1.0) {E};
        \node[tgt] (f) at ( 1.8,-1.0) {F};
        \foreach \n in {a,b,c,d,e,f}{\draw[->, dashed, UFSRed] (M) -- (\n);}
      \end{tikzpicture}\\
      \footnotesize Uma mudança obriga mexer em muitas classes.\\
      \textbf{Refatoração:} \emph{Move Method/Field}.
    \end{column}
  \end{columns}
\end{frame}

\begin{frame}[fragile]{Switch Statements / cadeias \texttt{if-elif}}
\begin{lstlisting}
def calcular_frete(pedido):
    if pedido.tipo == "NORMAL":   return pedido.peso * 5.0
    elif pedido.tipo == "EXPRESSO": return pedido.peso * 10.0
    elif pedido.tipo == "SEDEX":  return pedido.peso * 15.0
    else: raise ValueError(pedido.tipo)
\end{lstlisting}
  \footnotesize
  \textbf{Por que cheira mal?} adicionar um tipo novo exige caçar
  \emph{todos} os switches --- é \emph{Shotgun Surgery} esperando para
  acontecer.\\
  \textbf{Refatoração:} \emph{Replace Conditional with Polymorphism}
  (uma classe por tipo de envio).
\end{frame}

\begin{frame}{Recapitulando a Parte I}
  \begin{center}
  \begin{tabular}{ll}
    \toprule
    \textbf{Smell} & \textbf{Refatoração de combate} \\
    \midrule
    Long Method            & Extract Function \\
    Large Class            & Extract Class \\
    Long Parameter List    & Introduce Parameter Object \\
    Duplicated Code        & Extract Function / Pull Up \\
    Divergent Change       & Extract Class \\
    Shotgun Surgery        & Move Method / Field \\
    Switch Statements      & Replace Conditional w/ Polymorphism \\
    \bottomrule
  \end{tabular}
  \end{center}
\end{frame}

% =====================================================================
\part{Solução: Refatoração}
% =====================================================================
\section{O que é refatoração}

\begin{frame}{Refatoração: definição}
  \begin{block}{Definição (Fowler, 2018)}
    Refatorar é \emph{reestruturar código existente alterando sua estrutura
    interna sem mudar seu comportamento externo observável}.
  \end{block}
  \begin{exampleblock}{Intuição}
    Como reorganizar uma cozinha enquanto se cozinha: você muda onde as
    coisas ficam, mas o jantar continua sendo o mesmo --- e a cada passo
    você verifica que ainda está cozinhando.
  \end{exampleblock}
\end{frame}

\section{Catálogo de refatorações básicas}

\begin{frame}{Quatro refatorações que cobrem 80\% dos casos}
  \begin{center}
  \begin{tikzpicture}[
    box/.style={rectangle, draw=UFSBlue!60, fill=UFSBlue!10, rounded corners,
      minimum width=5.2cm, minimum height=1.2cm, font=\small,
      align=left, text width=5cm}
  ]
    \node[box] (a) at (-3.0, 1.2)
      {\textbf{Extract Function}\\
       \scriptsize trecho coeso $\to$ função com nome próprio};
    \node[box] (b) at ( 3.0, 1.2)
      {\textbf{Rename}\\
       \scriptsize trocar nome obscuro por nome que revela intenção};
    \node[box] (c) at (-3.0,-1.2)
      {\textbf{Inline}\\
       \scriptsize indireção desnecessária $\to$ código no lugar};
    \node[box] (d) at ( 3.0,-1.2)
      {\textbf{Introduce Parameter Object}\\
       \scriptsize grupo de parâmetros recorrentes $\to$ objeto};
  \end{tikzpicture}
  \end{center}
  \vfill
  \begin{center}
  \qrcode[height=1.8cm]{https://refactoring.guru}\quad
  \raisebox{0.6cm}{\small Catálogo completo: \texttt{refactoring.guru}}
  \end{center}
\end{frame}

\begin{frame}[fragile]{Extract Function --- antes}
\begin{lstlisting}
def imprimir_extrato(cliente, lancamentos):
    print(f"Cliente: {cliente.nome}")
    print("-" * 30)
    total = 0
    for l in lancamentos:
        total += l.valor
    print(f"Total: R$ {total:.2f}")
\end{lstlisting}
\small O cálculo do total está \emph{escondido} dentro de uma função de
impressão. O nome da função mente sobre o que ela faz.
\end{frame}

\begin{frame}[fragile]{Extract Function --- depois}
\begin{lstlisting}
def calcular_total(lancamentos):
    return sum(l.valor for l in lancamentos)

def imprimir_extrato(cliente, lancamentos):
    print(f"Cliente: {cliente.nome}")
    print("-" * 30)
    total = calcular_total(lancamentos)
    print(f"Total: R$ {total:.2f}")
\end{lstlisting}
\small Cada função tem \emph{um} motivo para mudar.
\texttt{calcular\_total} agora é testável isoladamente e reutilizável.
\end{frame}

\begin{frame}{Mecânica passo a passo de Extract Function}
  \begin{enumerate}
    \item Identifique um trecho coeso com responsabilidade própria.
    \item Crie uma função nova com nome que \emph{revele a intenção}.
    \item Copie o trecho para a nova função; declare parâmetros para
              cada variável usada de fora.
    \item Substitua o trecho original por uma chamada à nova função.
    \item \textbf{Rode os testes.} Se passarem, \emph{commit}.
  \end{enumerate}
  \vspace{6pt}
  {\small Passos pequenos $+$ testes a cada passo $=$ confiança.}
\end{frame}

\section{Quando refatorar}

\begin{frame}{Três gatilhos práticos}
  \begin{center}
  \begin{tikzpicture}[
    g/.style={rectangle, draw=green!50!black, fill=green!10,
      rounded corners, minimum width=10cm, minimum height=0.9cm,
      font=\small, align=left, text width=9.5cm}
  ]
    \node[g] (g1) at (0, 1.5)
      {\textbf{Regra dos três:} duplicou pela terceira vez? Extraia.};
    \node[g] (g2) at (0, 0.0)
      {\textbf{Antes de adicionar feature:} prepare o terreno onde
       a nova funcionalidade vai entrar (\emph{``make the change easy,
       then make the easy change''} -- Kent Beck).};
    \node[g] (g3) at (0,-1.5)
      {\textbf{Ao corrigir um bug:} se entender o bug foi difícil,
       refatore para que o próximo na mesma região seja óbvio.};
  \end{tikzpicture}
  \end{center}
\end{frame}

\begin{frame}{Quando \emph{não} refatorar}
  \begin{itemize}
    \item Código que está prestes a ser \emph{descartado}.
    \item Sem \emph{rede de testes} cobrindo o trecho --- escreva testes primeiro.
    \item Misturado com mudança de comportamento no mesmo \emph{commit}
          (separe: ou refatora, ou muda comportamento, nunca os dois juntos).
    \item Sob pressão extrema de prazo, sem combinar com o time.
  \end{itemize}
\end{frame}

\section{Estudo de caso}

\begin{frame}[fragile]{Ponto de partida (cheirando mal)}
\begin{lstlisting}
def calc(itens, t):
    s = 0
    for i in itens:
        s = s + i[0] * i[1]
    if t == "vip":
        s = s * 0.9
    elif t == "func":
        s = s * 0.8
    return s
\end{lstlisting}
\small Cheiros: nomes opacos (\texttt{s}, \texttt{i}, \texttt{t}),
tupla anônima como ``item'', desconto misturado com soma,
\texttt{if/elif} sobre \emph{strings mágicas}.
\end{frame}

\begin{frame}[fragile]{Depois das refatorações}
\begin{lstlisting}
DESCONTO = {"vip": 0.10, "func": 0.20}

def subtotal(itens):
    return sum(item.preco * item.qtd for item in itens)

def calcular_total(itens, perfil="comum"):
    return subtotal(itens) * (1 - DESCONTO.get(perfil, 0))
\end{lstlisting}
\small Aplicamos: \emph{Rename}, \emph{Extract Function}
(\texttt{subtotal}), \emph{Replace Conditional with Lookup}
(dicionário \texttt{DESCONTO}) e \emph{Introduce Parameter Object}
implícito (\texttt{item.preco} em vez de \texttt{i[0]}).
\end{frame}

% =====================================================================
\part{Disciplina: Ciclo, Testes e PEP 8}
% =====================================================================
\section{Ciclo Red--Green--Refactor}

\begin{frame}{Por que precisamos de um \emph{ciclo}?}
  \begin{columns}[T]
    \begin{column}{0.5\textwidth}
      \textbf{Problema:}\\[0.4em]
      Refatorar ``no instinto'' quebra código.
      Mudanças grandes misturam alteração de
      \emph{estrutura} com alteração de
      \emph{comportamento} --- e bugs se escondem.
    \end{column}
    \begin{column}{0.5\textwidth}
      \textbf{Solução:}\\[0.4em]
      Um \textbf{ciclo curto e disciplinado} em que
      cada passo é pequeno, verificável por testes e
      reversível. \emph{Red--Green--Refactor.}
    \end{column}
  \end{columns}
\end{frame}

\begin{frame}{O ciclo Red--Green--Refactor}
  \begin{center}
  \begin{tikzpicture}[
    state/.style={circle, draw, minimum size=1.8cm, font=\scriptsize\bfseries, align=center},
    arrow/.style={->, thick, >=Stealth}
  ]
    \node[state, fill=red!20]    (R) at (0,0)    {RED\\\tiny teste falha};
    \node[state, fill=green!25]  (G) at (3.6,0)  {GREEN\\\tiny teste passa};
    \node[state, fill=blue!15]   (F) at (1.8,-2.4) {REFACTOR\\\tiny limpe};

    \draw[arrow] (R) to[bend left=20] node[above, font=\tiny]{escreva o mínimo} (G);
    \draw[arrow] (G) to[bend left=20] node[right, font=\tiny]{melhore} (F);
    \draw[arrow] (F) to[bend left=20] node[left,  font=\tiny]{próximo teste} (R);
  \end{tikzpicture}
  \end{center}
  \begin{itemize}\footnotesize
    \item \textbf{Red:} teste que falha define o que o código deve fazer.
    \item \textbf{Green:} faça passar do jeito mais simples.
    \item \textbf{Refactor:} com a rede verde, melhore o design.
  \end{itemize}
\end{frame}

\begin{frame}{Os dois chapéus de Kent Beck}
  \begin{center}
  \begin{tikzpicture}[
    hat/.style={rectangle, rounded corners, draw, minimum width=5cm,
                minimum height=2.2cm, align=center, font=\small}
  ]
    \node[hat, fill=UFSYellow!40] (A) at (-3.5, 0)
      {\textbf{Chapéu 1: Adicionar Função}\\[0.3em]
       Mudar \emph{comportamento}.\\
       \emph{Não} mexer em estrutura.\\
       Testes novos podem falhar.};
    \node[hat, fill=UFSBlue!20] (B) at (3.5, 0)
      {\textbf{Chapéu 2: Refatorar}\\[0.3em]
       Mudar \emph{estrutura}.\\
       \emph{Não} mexer em comportamento.\\
       Testes \textbf{continuam} passando.};
    \draw[<->, thick, dashed] (A) -- (B)
      node[midway, above, font=\scriptsize\itshape]{nunca os dois ao mesmo tempo};
  \end{tikzpicture}
  \end{center}
\end{frame}

\begin{frame}[fragile]{Exemplo: um ciclo completo}
  \textbf{Red} --- teste primeiro:
\begin{lstlisting}
def test_total_com_desconto():
    assert total([10, 20, 30], desconto=0.1) == 54.0
\end{lstlisting}
  \textbf{Green} --- o mínimo que faz passar:
\begin{lstlisting}
def total(itens, desconto):
    return sum(itens) * (1 - desconto)
\end{lstlisting}
  \textbf{Refactor} --- nomes e responsabilidades:
\begin{lstlisting}
def _aplicar_desconto(valor, taxa):
    return valor * (1 - taxa)

def total(itens, desconto=0.0):
    return _aplicar_desconto(sum(itens), desconto)
\end{lstlisting}
\end{frame}

\section{Testes com \texttt{pytest}}

\begin{frame}{\texttt{pytest}: definição}
  \begin{block}{Definição}
    \texttt{pytest} é um framework de testes que descobre
    automaticamente funções \texttt{test\_*} em arquivos \texttt{test\_*.py},
    executa-as, e relata falhas com introspecção detalhada da expressão \texttt{assert}.
  \end{block}
  \begin{exampleblock}{Intuição}
    Você escreve Python comum (\texttt{def} + \texttt{assert});
    o framework cuida da descoberta, do relatório e da
    rede de segurança que sustenta o ciclo Red--Green--Refactor.
  \end{exampleblock}
\end{frame}

\begin{frame}[fragile]{Anatomia de um teste}
\begin{lstlisting}
# arquivo: test_carrinho.py
from carrinho import total

def test_total_lista_vazia():
    assert total([]) == 0

def test_total_soma_itens():
    assert total([1, 2, 3]) == 6
\end{lstlisting}
Execução:
\begin{lstlisting}[language=bash, keywordstyle=]
$ pytest -v
test_carrinho.py::test_total_lista_vazia   PASSED
test_carrinho.py::test_total_soma_itens    PASSED
\end{lstlisting}
\end{frame}

\begin{frame}{Recursos essenciais do \texttt{pytest}}
  \begin{itemize}
    \item \textbf{Fixtures} (\texttt{@pytest.fixture}) --- preparam dados/estado reutilizáveis.
    \item \textbf{Parametrização} (\texttt{@pytest.mark.parametrize}) --- um teste, vários casos.
    \item \textbf{Exceções esperadas} (\texttt{with pytest.raises(...)}).
    \item \textbf{Descoberta automática} --- arquivos \texttt{test\_*.py}, funções \texttt{test\_*}.
  \end{itemize}
\end{frame}

\begin{frame}[fragile]{Parametrização e exceções}
\begin{lstlisting}
import pytest

@pytest.mark.parametrize("entrada,esperado", [
    ([],       0),
    ([1, 2],   3),
    ([10]*5, 50),
])
def test_total(entrada, esperado):
    assert total(entrada) == esperado

def test_sacar_mais_que_saldo_falha():
    c = Conta(saldo=100)
    with pytest.raises(SaldoInsuficienteError):
        c.sacar(150)
\end{lstlisting}
\small Um teste parametrizado roda várias vezes (uma falha não impede as outras);
\texttt{pytest.raises} \emph{documenta} o contrato de exceções.
\end{frame}

\section{Estilo: PEP 8}

\begin{frame}{PEP 8: definição}
  \begin{block}{Definição}
    \textbf{PEP 8} (\emph{Python Enhancement Proposal 8}) é o documento
    que padroniza a aparência do código Python: indentação,
    nomenclatura, comprimento de linha, espaçamento em torno de
    operadores, organização de \texttt{import}s, entre outros.
  \end{block}
  \begin{exampleblock}{Intuição --- ``a hobgoblin of little minds''}
    Consistência importa, mas não é cega: se a PEP 8 atrapalha
    a legibilidade num caso específico, ignore-a \emph{nesse} caso.
    Consistência local $>$ consistência global $>$ PEP 8 literal.
  \end{exampleblock}
\end{frame}

\begin{frame}{Regras essenciais da PEP 8}
  \begin{center}
  \begin{tikzpicture}[
    rule/.style={rectangle, draw, rounded corners, fill=UFSBlue!8,
                 minimum width=11cm, minimum height=0.7cm,
                 font=\small, align=left, inner sep=6pt, text width=10.6cm}
  ]
    \node[rule] (r1) at (0, 2.4) {\textbf{Indentação:} 4 espaços por nível (nunca tab).};
    \node[rule] (r2) at (0, 1.6) {\textbf{Linha:} máx.\ 79 caracteres de código, 72 de comentário.};
    \node[rule] (r3) at (0, 0.8) {\textbf{Nomes:} \texttt{snake\_case} (funções), \texttt{PascalCase} (classes), \texttt{MAIUSCULAS} (constantes).};
    \node[rule] (r4) at (0, 0.0) {\textbf{Imports:} um por linha, no topo, agrupados (stdlib, terceiros, locais).};
    \node[rule] (r5) at (0,-0.8) {\textbf{Espaçamento:} 1 espaço em torno de \texttt{=}, exceto em argumentos default.};
    \node[rule] (r6) at (0,-1.6) {\textbf{Comparações:} \texttt{is None} (não \texttt{== None}); \texttt{if x:} para coleções vazias.};
  \end{tikzpicture}
  \end{center}
\end{frame}

\begin{frame}[fragile]{Antes e depois da PEP 8}
\textbf{Antes} --- viola convenções:
\begin{lstlisting}
import os,sys
def CalcularMedia( numeros ):
  Total=0
  for N in numeros:Total=Total+N
  return Total/len( numeros )
\end{lstlisting}
\textbf{Depois} --- de acordo com a PEP 8:
\begin{lstlisting}
import os
import sys


def calcular_media(numeros):
    return sum(numeros) / len(numeros)
\end{lstlisting}
\footnotesize\emph{Mesmo comportamento, leitura imediata.} É uma refatoração:
estrutura mudou, comportamento não.
\end{frame}

\begin{frame}[fragile]{Ferramentas que aplicam PEP 8 automaticamente}
  \begin{itemize}
    \item \texttt{ruff} --- linter rapidíssimo, agrupa \texttt{flake8} e mais.
    \item \texttt{black} --- formatador opinativo, padroniza tudo.
    \item \texttt{isort} --- ordena \texttt{import}s.
  \end{itemize}
  \vspace{0.5em}
\begin{lstlisting}[language=bash, keywordstyle=]
$ ruff check src/
$ black src/
$ pytest
\end{lstlisting}
  \footnotesize\emph{Pipeline típico:} formatar $\to$ lintar $\to$ testar.
\end{frame}

\section{Integrando os três}

\begin{frame}{O fluxo de trabalho profissional}
  \begin{center}
  \begin{tikzpicture}[
    node distance=1.0cm and 1.2cm,
    box/.style={rectangle, rounded corners, draw, minimum width=3cm,
                minimum height=1cm, align=center, font=\small},
    arrow/.style={->, thick, >=Stealth}
  ]
    \node[box, fill=UFSRed!25]    (red)   {Escreva teste\\(\texttt{pytest})};
    \node[box, fill=green!30,  right=of red]   (green) {Implemente\\(funciona)};
    \node[box, fill=UFSBlue!20, right=of green] (refac) {Refatore\\estrutura};
    \node[box, fill=UFSYellow!40, below=of green] (style) {Formate\\(\texttt{ruff}/\texttt{black})};

    \draw[arrow] (red)   -- (green);
    \draw[arrow] (green) -- (refac);
    \draw[arrow] (refac) -- (style);
    \draw[arrow] (style.west) to[bend left=25] (red.south);
  \end{tikzpicture}
  \end{center}
  \begin{itemize}\footnotesize
    \item Cheiros (Parte I) disparam refatorações nomeadas (Parte II).
    \item Testes garantem que a refatoração não muda comportamento.
    \item PEP 8 garante que o código refatorado é legível pela equipe.
  \end{itemize}
\end{frame}

\begin{frame}{Encerramento: próximos passos}
  \begin{columns}[T]
    \begin{column}{0.5\textwidth}
      \begin{block}{Pratique}
        \begin{itemize}\small
          \item Identifique 3 cheiros em um projeto seu.
          \item Cubra-os com testes \texttt{pytest} antes de mexer.
          \item Aplique a refatoração de combate em passos pequenos.
          \item Rode \texttt{ruff} + \texttt{black} ao final.
        \end{itemize}
      \end{block}
    \end{column}
    \begin{column}{0.5\textwidth}
      \begin{exampleblock}{Leituras recomendadas}
        \begin{itemize}\small
          \item Fowler (2018), \emph{Refactoring}, caps.\ 1 e 3.
          \item Beck (2003), \emph{TDD by Example}, parte I.
          \item Martin (2008), \emph{Clean Code}.
          \item PEP 8 --- \texttt{peps.python.org/pep-0008}.
        \end{itemize}
      \end{exampleblock}
    \end{column}
  \end{columns}
  \vspace{0.5em}
  \begin{center}
    \emph{``Make it work, make it right, make it fast.''} --- Kent Beck
  \end{center}
\end{frame}

% =====================================================================
\section*{Referências}
\begin{frame}[allowframebreaks]{Referências}
  \nocite{*}
  \bibliographystyle{plain}
  \bibliography{referencias}
\end{frame}

\end{document}
